% Options for packages loaded elsewhere
% Options for packages loaded elsewhere
\PassOptionsToPackage{unicode}{hyperref}
\PassOptionsToPackage{hyphens}{url}
\PassOptionsToPackage{dvipsnames,svgnames,x11names}{xcolor}
%
\documentclass[
  letterpaper,
  DIV=11,
  numbers=noendperiod]{scrreprt}
\usepackage{xcolor}
\usepackage{amsmath,amssymb}
\setcounter{secnumdepth}{5}
\usepackage{iftex}
\ifPDFTeX
  \usepackage[T1]{fontenc}
  \usepackage[utf8]{inputenc}
  \usepackage{textcomp} % provide euro and other symbols
\else % if luatex or xetex
  \usepackage{unicode-math} % this also loads fontspec
  \defaultfontfeatures{Scale=MatchLowercase}
  \defaultfontfeatures[\rmfamily]{Ligatures=TeX,Scale=1}
\fi
\usepackage{lmodern}
\ifPDFTeX\else
  % xetex/luatex font selection
\fi
% Use upquote if available, for straight quotes in verbatim environments
\IfFileExists{upquote.sty}{\usepackage{upquote}}{}
\IfFileExists{microtype.sty}{% use microtype if available
  \usepackage[]{microtype}
  \UseMicrotypeSet[protrusion]{basicmath} % disable protrusion for tt fonts
}{}
\makeatletter
\@ifundefined{KOMAClassName}{% if non-KOMA class
  \IfFileExists{parskip.sty}{%
    \usepackage{parskip}
  }{% else
    \setlength{\parindent}{0pt}
    \setlength{\parskip}{6pt plus 2pt minus 1pt}}
}{% if KOMA class
  \KOMAoptions{parskip=half}}
\makeatother
% Make \paragraph and \subparagraph free-standing
\makeatletter
\ifx\paragraph\undefined\else
  \let\oldparagraph\paragraph
  \renewcommand{\paragraph}{
    \@ifstar
      \xxxParagraphStar
      \xxxParagraphNoStar
  }
  \newcommand{\xxxParagraphStar}[1]{\oldparagraph*{#1}\mbox{}}
  \newcommand{\xxxParagraphNoStar}[1]{\oldparagraph{#1}\mbox{}}
\fi
\ifx\subparagraph\undefined\else
  \let\oldsubparagraph\subparagraph
  \renewcommand{\subparagraph}{
    \@ifstar
      \xxxSubParagraphStar
      \xxxSubParagraphNoStar
  }
  \newcommand{\xxxSubParagraphStar}[1]{\oldsubparagraph*{#1}\mbox{}}
  \newcommand{\xxxSubParagraphNoStar}[1]{\oldsubparagraph{#1}\mbox{}}
\fi
\makeatother


\usepackage{longtable,booktabs,array}
\newcounter{none} % for unnumbered tables
\usepackage{calc} % for calculating minipage widths
% Correct order of tables after \paragraph or \subparagraph
\usepackage{etoolbox}
\makeatletter
\patchcmd\longtable{\par}{\if@noskipsec\mbox{}\fi\par}{}{}
\makeatother
% Allow footnotes in longtable head/foot
\IfFileExists{footnotehyper.sty}{\usepackage{footnotehyper}}{\usepackage{footnote}}
\makesavenoteenv{longtable}
\usepackage{graphicx}
\makeatletter
\newsavebox\pandoc@box
\newcommand*\pandocbounded[1]{% scales image to fit in text height/width
  \sbox\pandoc@box{#1}%
  \Gscale@div\@tempa{\textheight}{\dimexpr\ht\pandoc@box+\dp\pandoc@box\relax}%
  \Gscale@div\@tempb{\linewidth}{\wd\pandoc@box}%
  \ifdim\@tempb\p@<\@tempa\p@\let\@tempa\@tempb\fi% select the smaller of both
  \ifdim\@tempa\p@<\p@\scalebox{\@tempa}{\usebox\pandoc@box}%
  \else\usebox{\pandoc@box}%
  \fi%
}
% Set default figure placement to htbp
\def\fps@figure{htbp}
\makeatother





\setlength{\emergencystretch}{3em} % prevent overfull lines

\providecommand{\tightlist}{%
  \setlength{\itemsep}{0pt}\setlength{\parskip}{0pt}}



 


\KOMAoption{captions}{tableheading}
\makeatletter
\@ifpackageloaded{bookmark}{}{\usepackage{bookmark}}
\makeatother
\makeatletter
\@ifpackageloaded{caption}{}{\usepackage{caption}}
\AtBeginDocument{%
\ifdefined\contentsname
  \renewcommand*\contentsname{Table of contents}
\else
  \newcommand\contentsname{Table of contents}
\fi
\ifdefined\listfigurename
  \renewcommand*\listfigurename{List of Figures}
\else
  \newcommand\listfigurename{List of Figures}
\fi
\ifdefined\listtablename
  \renewcommand*\listtablename{List of Tables}
\else
  \newcommand\listtablename{List of Tables}
\fi
\ifdefined\figurename
  \renewcommand*\figurename{Figure}
\else
  \newcommand\figurename{Figure}
\fi
\ifdefined\tablename
  \renewcommand*\tablename{Table}
\else
  \newcommand\tablename{Table}
\fi
}
\@ifpackageloaded{float}{}{\usepackage{float}}
\floatstyle{ruled}
\@ifundefined{c@chapter}{\newfloat{codelisting}{h}{lop}}{\newfloat{codelisting}{h}{lop}[chapter]}
\floatname{codelisting}{Listing}
\newcommand*\listoflistings{\listof{codelisting}{List of Listings}}
\makeatother
\makeatletter
\makeatother
\makeatletter
\@ifpackageloaded{caption}{}{\usepackage{caption}}
\@ifpackageloaded{subcaption}{}{\usepackage{subcaption}}
\makeatother
\usepackage{bookmark}
\IfFileExists{xurl.sty}{\usepackage{xurl}}{} % add URL line breaks if available
\urlstyle{same}
\hypersetup{
  pdftitle={Ecosystem Service Hotspots: A Global Policy Analysis},
  pdfauthor={Jerónimo Rodríguez Escobar, Becky Chaplin-Kramer, Richard P. Sharp, [Other Authors]},
  colorlinks=true,
  linkcolor={blue},
  filecolor={Maroon},
  citecolor={Blue},
  urlcolor={Blue},
  pdfcreator={LaTeX via pandoc}}


\title{Ecosystem Service Hotspots: A Global Policy Analysis}
\usepackage{etoolbox}
\makeatletter
\providecommand{\subtitle}[1]{% add subtitle to \maketitle
  \apptocmd{\@title}{\par {\large #1 \par}}{}{}
}
\makeatother
\subtitle{Where nature is changing, who is affected, and what drives
change}
\author{Jerónimo Rodríguez Escobar, Becky Chaplin-Kramer, Richard P. Sharp, [Other Authors]}
\date{2026-06-05}
\begin{document}
\maketitle

\renewcommand*\contentsname{Table of contents}
{
\setcounter{tocdepth}{2}
\tableofcontents
}

\bookmarksetup{startatroot}

\chapter*{Preface}\label{preface}
\addcontentsline{toc}{chapter}{Preface}

\markboth{Preface}{Preface}

\bookmarksetup{startatroot}

\chapter*{Executive Summary}\label{executive-summary}
\addcontentsline{toc}{chapter}{Executive Summary}

\markboth{Executive Summary}{Executive Summary}

This workflow brings together global data on ecosystem services, land
cover, and people to help us understand where nature is changing, who is
affected, and where action is most needed. It combines data processing,
change detection, and hotspot mapping in a way that is transparent and
reproducible.

\textbf{Why does it matter?} By identifying areas of rapid change or
high importance, this pipeline supports better decision-making for
conservation, policy, and sustainable development.

\section*{Glossary}\label{glossary}
\addcontentsline{toc}{section}{Glossary}

\markright{Glossary}

\begin{itemize}
\tightlist
\item
  \textbf{AOO}: Area of Occupancy. A standard 10 km equal-area grid used
  for spatial analysis.
\item
  \textbf{ES}: Ecosystem Services. Benefits people obtain from nature
  (e.g., pollination, coastal protection).
\item
  \textbf{Hotspot}: A grid cell showing unusually high or low relative
  change in ecosystem services (the extreme 5\% tail of the
  distribution).
\item
  \textbf{KS Analysis}: Kolmogorov-Smirnov test, a statistical method
  used here to compare the socioeconomic profiles of hotspots
  vs.~non-hotspots.
\item
  \textbf{Zonal Statistics}: Calculations that summarize high-resolution
  raster data within the boundaries of polygons (grid cells or regions).
\end{itemize}

\bookmarksetup{startatroot}

\chapter*{Key Takeaways}\label{key-takeaways}
\addcontentsline{toc}{chapter}{Key Takeaways}

\markboth{Key Takeaways}{Key Takeaways}

\subsection*{1. What is a ``Hotspot''? (Methodological
Context)}\label{what-is-a-hotspot-methodological-context}
\addcontentsline{toc}{subsection}{1. What is a ``Hotspot''?
(Methodological Context)}

Before interpreting the results, it is critical to understand how a
hotspot is defined in this analysis: * \textbf{A Relative Extreme, Not
an Absolute Threshold:} A hotspot is a grid cell that falls in the most
extreme \textbf{5\% of relative change} (using Symmetric Percentage
Change) globally for a specific service. * \textbf{Intensity Over
Volume:} We use percentage change rather than absolute volume to
normalize for the size of the baseline ecosystem. * \textbf{Focus on
Decline:} For this phase of the analysis, hotspots represent
\textbf{worry spots}---areas experiencing the steepest declines in
service provision or the sharpest increases in environmental
risk/damage.

\subsection*{2. Geographic Clustering \& Compound
Risk}\label{geographic-clustering-compound-risk}
\addcontentsline{toc}{subsection}{2. Geographic Clustering \& Compound
Risk}

Hotspots of ecosystem service decline are not randomly distributed; they
cluster heavily in the Global South, creating epicenters of
\textbf{Compound Risk} where communities face multiple simultaneous
ecological failures. * \textbf{The Epicenters of Hotness:} Latin America
\& the Caribbean and East Asia \& the Pacific are the most intensely
affected regions globally. * \textbf{Disproportionate Burden:} Latin
America and Sub-Saharan Africa carry a vastly disproportionate share of
the global damage compared to their landmass.

\subsection*{3. Socioeconomic Profiling (Who is
affected?)}\label{socioeconomic-profiling-who-is-affected}
\addcontentsline{toc}{subsection}{3. Socioeconomic Profiling (Who is
affected?)}

Using Kolmogorov-Smirnov (KS) tests, we compared the socioeconomic
profile of hotspots against the ``business-as-usual'' median landscape.
* \textbf{The Urbanization / Wealth Correlation:} Declines in Nature
Access, Coastal Risk Reduction, and Retention Ratios consistently occur
in grid cells with higher Population Density, higher Built Area, and
higher GDP than the global median. * \textbf{The Rural Divide:}
Conversely, hotspots of Pollination loss and Sediment Export increases
are uniquely concentrated in less urbanized, lower-density, and highly
agricultural contexts.

\subsection*{4. Land Cover Drivers (The ``Attribution
Gap'')}\label{land-cover-drivers-the-attribution-gap}
\addcontentsline{toc}{subsection}{4. Land Cover Drivers (The
``Attribution Gap'')}

To move beyond correlation, we assessed the direct overlap between ES
hotspots and the top 5\% of physical Land Cover Conversion. *
\textbf{The Attribution Gap:} A significant portion of ecosystem service
hotspots do not overlap with major land conversion hotspots within the
same 10km cell, pointing to the critical role of degradation or upstream
drivers.

\subsection*{5. Absolute Population
Exposure}\label{absolute-population-exposure}
\addcontentsline{toc}{subsection}{5. Absolute Population Exposure}

While socioeconomic profiling reveals relative characteristics, absolute
population provides the scale of human impact. * \textbf{Concentration
in Middle-Income Nations:} The vast majority of people exposed to ES
decline hotspots reside in Upper and Lower Middle-Income countries.

\part{Part 1: Understanding the Challenge}

\chapter{The Challenge}\label{the-challenge}

Why ecosystem services matter and where we stand

\hfill\break

\section{Global Ecosystem Services Under
Pressure}\label{global-ecosystem-services-under-pressure}

The well-being of billions of people is intrinsically linked to the
health of ecosystems. These life-support systems, often called
\textbf{ecosystem services}, include:

\begin{itemize}
\tightlist
\item
  \textbf{Water purification and erosion control} -- Nitrogen and
  sediment retention
\item
  \textbf{Food production} -- Pollination of crops; fisheries
  productivity
\item
  \textbf{Climate regulation} -- Carbon sequestration; temperature
  buffering
\item
  \textbf{Coastal protection} -- Mangroves and coral reefs reducing
  storm surge impact
\item
  \textbf{Cultural value} -- Recreation, spiritual significance of
  natural areas
\end{itemize}

Yet these life-support systems are \textbf{changing rapidly and
unevenly} across the globe. Between 1992 and 2020:

\begin{itemize}
\tightlist
\item
  Forest cover declined by \textasciitilde X million hectares globally
\item
  Agricultural intensification accelerated in developing regions
\item
  Wetland conversion for urban/agricultural use reached critical levels
\item
  Climate change intensified rainfall variability and droughts
\end{itemize}

\textbf{The problem}: Ecosystem service decline is
\textbf{geographically uneven}. Some regions face catastrophic loss
while others remain relatively stable. Worse, the \textbf{burden falls
disproportionately on low-income countries and vulnerable populations}
least able to adapt.

\section{Decision-Makers Face an Information
Gap}\label{decision-makers-face-an-information-gap}

Conservation and development professionals operate in a sea of data but
often lack answers to three critical questions:

\subsection{\texorpdfstring{🌍 \textbf{WHERE are ecosystem services
declining
most?}}{🌍 WHERE are ecosystem services declining most?}}\label{where-are-ecosystem-services-declining-most}

Global datasets exist, but decision-makers need geographically explicit,
actionable priorities. Where should conservation investment flow? Which
landscapes require immediate intervention?

\subsection{\texorpdfstring{👥 \textbf{WHO is
affected?}}{👥 WHO is affected?}}\label{who-is-affected}

Not all ecosystem service loss is equally consequential. Loss in
low-income agricultural zones threatens food security; loss in urban
regions threatens water supply. Understanding affected populations is
crucial for equitable, effective policy.

\subsection{\texorpdfstring{🔧 \textbf{WHY are services
declining?}}{🔧 WHY are services declining?}}\label{why-are-services-declining}

Is decline driven by land conversion (deforestation, urban sprawl) or
land degradation (quality declining within remaining natural areas)? The
answer determines the intervention strategy: - \textbf{Conversion-driven
decline} → Protect remaining land, strengthen tenure rights -
\textbf{Degradation-driven decline} → Improve land management, restore
ecosystem function

\section{Our Approach: A Three-Question
Framework}\label{our-approach-a-three-question-framework}

This book integrates global ecosystem service models, land cover data,
and socioeconomic layers to answer WHERE/WHO/WHY systematically:

\begin{enumerate}
\def\labelenumi{\arabic{enumi}.}
\tightlist
\item
  \textbf{WHERE Analysis} (Chapters 4-5)

  \begin{itemize}
  \tightlist
  \item
    Identified hotspots: regions in the extreme top/bottom 5\% of
    ecosystem service change
  \item
    Geographic patterns: which regions, income groups, and biomes bear
    disproportionate burden
  \item
    Compound risk: areas where multiple services fail simultaneously
  \end{itemize}
\item
  \textbf{WHO Analysis} (Chapter 6)

  \begin{itemize}
  \tightlist
  \item
    Population exposure: how many people live in hotspot regions
  \item
    Socioeconomic profiles: are hotspot populations wealthy, poor,
    rural, urban?
  \item
    Statistical validation: are hotspot populations truly different from
    global norm
  \end{itemize}
\item
  \textbf{WHY Analysis} (Chapter 7)

  \begin{itemize}
  \tightlist
  \item
    Attribution: what \% of decline coincides with observed land cover
    change
  \item
    Attribution gap: what portion of decline is unexplained (degradation
    signal)
  \item
    Service-specific drivers: which mechanisms drive which services
  \end{itemize}
\end{enumerate}

\section{Why This Analysis Matters}\label{why-this-analysis-matters}

\textbf{For conservation organizations}: Direct limited resources to
regions of greatest need, vulnerable populations

\textbf{For governments}: Identify priorities for environmental policy,
land-use planning, climate adaptation strategies

\textbf{For international donors}: Allocate development aid
strategically, considering both ecosystem and equity dimensions

\textbf{For research}: Understand spatial and mechanistic patterns in
ecosystem service change; ground-truth hypotheses

\textbf{For corporations}: Manage supply chain risks; invest in
nature-positive initiatives where impacts are greatest

\section{How to Use This Book}\label{how-to-use-this-book}

\textbf{Quick read (1-2 hours)}: - Chapter 1: Problem statement (this
chapter) - Chapter 2: Methods overview - Chapter 3: Global summary -
Chapter 8: Conclusions \& implications

\textbf{Policy-focused deep dive (2-3 hours)}: - Chapters 1-3 + 8
(above) - Then skip to Chapter 7 (Regional profiles) for your geography
of interest

\textbf{Complete analysis (4-5 hours)}: - Read all chapters in order -
Explore interactive tables and figures - Consult Technical Appendix for
methodological details

\textbf{For researchers}: - Chapters 2 (methods), 6 (WHO/socioeconomic),
7 (WHY/drivers) - Technical Appendix for full methodological detail and
validation

\section{Caveats \& Transparency}\label{caveats-transparency}

This analysis synthesizes global data at 10-km resolution. Important
limitations:

\begin{itemize}
\tightlist
\item
  \textbf{Coarse resolution} -- Sub-grid heterogeneity not captured;
  local variation substantial
\item
  \textbf{Data quality varies} -- Ecosystem service models depend on
  input data; validation is ongoing
\item
  \textbf{Temporal snapshot} -- 1992-2020 comparison; recent trends may
  differ
\item
  \textbf{Attribution gap} -- \textasciitilde30-55\% of decline
  unexplained; reflects knowledge gap not data gap
\item
  \textbf{Natural variation} -- Climate variability contributes to
  observed changes; attribution to human drivers incomplete
\end{itemize}

We embrace these uncertainties. This analysis is a starting point for
prioritization and hypothesis-driven research, not a definitive answer.

\section{About the Analysis}\label{about-the-analysis}

\textbf{Authors}: Jerónimo Rodríguez Escobar, Becky Chaplin-Kramer, Richard P. Sharp, [Other Authors]

\textbf{Data sources}:
\begin{itemize}
\tightlist
\item
  InVEST ecosystem service models (Spring/Natural Capital Project methods; \url{https://essd.copernicus.org/articles/17/3473/2025/})
\item
  ESA Climate Change Initiative land cover
\item
  IUCN 10km equal-area grids: Murray, Nicholas (2017). Global 10 x 10-km grids suitable for use in IUCN Red List of Ecosystems assessments (vector and raster format). figshare. Dataset. \url{https://doi.org/10.6084/m9.figshare.4653439.v1}
\item
  Socioeconomic and demographic data: World Bank regions; Global Population Density grids; UN Human Development Index (HDI) (Sherman, L., Proctor, J., Druckenmiller, H., Tapia, H., \& Hsiang, S. (2026). Global high-resolution estimates of the UN Human Development Index using satellite imagery and machine learning. Nature Communications, 17(1), 1315. \url{https://doi.org/10.1038/s41467-026-68805-6})
\end{itemize}

\textbf{Methods}: Python-based zonal statistics → R processing → hotspot
detection (top/bottom 5\% by Symmetric Percentage Change) →
Kolmogorov-Smirnov statistical profiling

\textbf{Reproducibility}: All code and intermediate outputs available;
see Technical Appendix

\section{What's Next}\label{whats-next}

Chapter 2 walks through the simplified methodology. Then we dive into
geographic (WHERE), socioeconomic (WHO), and mechanistic (WHY) analyses.

\chapter{Our Approach}\label{our-approach}

How we detected and analyzed ecosystem service hotspots

\hfill\break

\section{Data Sources}\label{data-sources}

This analysis integrates four major datasets:

\subsection{1. Ecosystem Service Models
(InVEST)}\label{ecosystem-service-models-invest}

\textbf{Source}: The Nature Conservancy's InVEST model suite
\textbf{Coverage}: Global, 300m resolution, 1992 \& 2020 snapshots
\textbf{Services modeled}: - Pollination (agricultural crop pollination)
- Nitrogen Export \& Retention (water quality) - Sediment Export \&
Retention (water quality) - Nature Access (recreation/cultural ecosystem
services) - Coastal Protection (mangrove/coral reef storm buffering)

\textbf{How it works}: InVEST uses land use, climate, and biophysical
data to estimate ecosystem service supply. A full technical description
is in the Technical Appendix.

\subsection{2. Land Cover Data}\label{land-cover-data}

\textbf{Source}: ESA Climate Change Initiative, 300m resolution
\textbf{Classification}: Reclassified to Transformed (agriculture,
urban, plantations) vs.~Natural (forests, grasslands, wetlands)
\textbf{Purpose}: Explain ecosystem service change; attribute decline to
land conversion vs.~degradation

\subsection{3. Administrative Grid}\label{administrative-grid}

\textbf{Source}: IUCN Oecism Opportunity Grid, 10km equal-area cells
\textbf{Coverage}: \textasciitilde1.5M cells covering all land globally
\textbf{Attributes}: Country, World Bank region, income group, WWF biome
\textbf{Purpose}: Common spatial unit for all analysis; enables
regional/income aggregation

\subsection{4. Socioeconomic Layers}\label{socioeconomic-layers}

\textbf{Sources}: World Bank, UNDP, CIESIN \textbf{Variables}:
Population density, GDP per capita, Human Development Index, built area,
agricultural intensity, Gini coefficient \textbf{Purpose}: Profile who
lives in hotspot regions; test whether hotspot populations differ from
global norm

\section{The Two-Path Analysis}\label{the-two-path-analysis}

We implement a \textbf{two-path analysis} to separate WHERE hotspots are
from WHY they occur:

\begin{verbatim}
┌─────────────────────────────────────────────────────────────┐
│ Path A: Global Trajectories (Pixel-level summary)           │
│ What: Aggregate 300m InVEST rasters to understand          │
│ global-level change trends by service                       │
│ Output: Global summary statistics per service               │
│ Use: Chapter 3 (global context)                             │
└─────────────────────────────────────────────────────────────┘

┌─────────────────────────────────────────────────────────────┐
│ Path B: Grid Analysis (Canonical for hotspots)              │
│ What: Zonal statistics within IUCN 10km cells               │
│ (difference of aggregates, not aggregate of differences)    │
│ Process:                                                     │
│   1. Aggregate 300m InVEST to 10km cell level (1992 & 2020)│
│   2. Calculate change per cell (absolute & percentage)      │
│   3. Identify hotspots: top/bottom 5% of distribution       │
│   4. Aggregate hotspots to regions/countries/biomes         │
│ Output: Hotspot locations, intensity, populations affected  │
│ Use: Chapters 4-7 (geographic, socioeconomic, drivers)      │
└─────────────────────────────────────────────────────────────┘
\end{verbatim}

\textbf{Why two paths?} Path A answers ``what happened globally''; Path
B answers ``where did it happen'' with actionable spatial detail.

\section{Change Metrics: From Absolute to Symmetric Percentage
Change}\label{change-metrics-from-absolute-to-symmetric-percentage-change}

\subsection{Absolute Change}\label{absolute-change}

\textbf{Formula}: \(\Delta = T_{2020} - T_{1992}\)

Simple subtraction. Problem: Services with different baselines are not
comparable (1 unit drop in service with baseline 10 ≠ 1 unit drop in
service with baseline 1000).

\subsection{Percentage Change (Naive)}\label{percentage-change-naive}

\textbf{Formula}:
\(\text{PC} = \frac{T_{2020} - T_{1992}}{T_{1992}} \times 100\%\)

Normalizes by baseline. Problem: \textbf{Zero-baseline services}
(services that could have 0 value in 1992) produce infinite values.

\subsection{\texorpdfstring{Symmetric Percentage Change (SPC) ✓
\emph{Our
choice}}{Symmetric Percentage Change (SPC) ✓ Our choice}}\label{symmetric-percentage-change-spc-our-choice}

\textbf{Formula}:
\(\text{SPC} = \frac{T_{2020} - T_{1992}}{(|T_{2020}| + |T_{1992}|) / 2} \times 100\%\)

Handles zero-baselines by using the mean of absolute values as
denominator. Symmetric treatment of gains/losses.

\textbf{Example}: - Service goes from 100 → 200: SPC = +67\% (2/3
increase) - Service goes from 100 → 0: SPC = -100\% (complete loss) -
Service goes from 0 → 100: SPC = +100\% (from nothing to something)

\textbf{Why SPC?} It prevents scale bias (large-baseline services
dominating hotspot identification) and handles zero-baseline services
elegantly.

\section{Hotspot Definition}\label{hotspot-definition}

\textbf{Hotspots} = Grid cells in the \textbf{top or bottom 5\%} of the
SPC distribution for each service.

\textbf{Rationale}: - Top 5\% = Locations with \textbf{most rapid
DECLINE} (greatest concern) - Bottom 5\% = Locations with \textbf{most
rapid RECOVERY} (positive learning) - 5\% threshold: Balance between
statistical rigor and practical granularity

\textbf{Result}: \textasciitilde75,000 hotspot cells per service (5\% of
1.5M global cells)

\section{Socioeconomic Profiling: Kolmogorov-Smirnov
Tests}\label{socioeconomic-profiling-kolmogorov-smirnov-tests}

We ask: \textbf{Are populations in hotspots statistically different from
non-hotspot populations?}

\textbf{Method}: Kolmogorov-Smirnov (KS) test compares cumulative
distributions: - \textbf{Hotspot population}: All socioeconomic values
in hotspot cells - \textbf{Non-hotspot population}: All values in
non-hotspot cells

\textbf{Output}: - \textbf{KS Statistic (D)}: Maximum vertical distance
between curves (0-1 scale); larger = more different - \textbf{P-value}:
Probability difference occurred by chance; p \textless{} 0.05 =
significant - \textbf{Cliff's Delta}: Effect size (how meaningful is the
difference, not just statistically significant)

\textbf{Interpretation}: - Significant KS test + high Cliff's Delta =
Hotspot populations meaningfully different (e.g., ``Pollination hotspots
are predominantly rural/agricultural'') - Significant but low Cliff's
Delta = Difference is small despite statistical significance

\section{Attribution Analysis: Overlap with Land Cover
Change}\label{attribution-analysis-overlap-with-land-cover-change}

\textbf{Question}: What \% of ecosystem service hotspots overlap with
observed land cover change?

\textbf{Method}: 1. Identify all grid cells with converted land cover
(natural → transformed, 1992→2020) 2. For each hotspot cell, check if it
also has conversion 3. Calculate \% overlap

\textbf{Result}: Heatmap showing which services' hotspots coincide with
which LCC types (forest loss, cropland expansion, etc.)

\textbf{Attribution Gap}: = 100\% - (\% overlap with any LCC type) -
High gap (e.g., 50\%) → Service decline not fully explained by observed
land conversion; suggests degradation signal - Low gap (e.g., 10\%) →
Service decline primarily explained by land cover change

\section{Quality Assurance}\label{quality-assurance}

\textbf{Validation steps}: 1. Grid geometry verification (all
\textasciitilde1.5M cells registered correctly) 2. Data integrity checks
(no missing values in key attributes) 3. Sanity checks on InVEST values
(within expected ranges; realistic spatial patterns) 4. Validation with
literature (hotspot locations match known conservation priorities)

\textbf{Uncertainties acknowledged}: - InVEST models depend on input
data quality (see Technical Appendix) - 10-km resolution may mask
sub-grid heterogeneity - 1992-2020 is a single temporal snapshot;
intermediate dynamics unknown - Socioeconomic data at 10-km resolution
may not capture local variation

\section{Reproducibility}\label{reproducibility}

All analysis is reproducible: - \textbf{Code}: R/Quarto notebooks
(available on project repository) - \textbf{Data}: Processed CSVs and
intermediate outputs in \texttt{/data/processed/} - \textbf{Full
pipeline}: Python zonal statistics → R processing → hotspot detection -
\textbf{See}: Technical Appendix for full methodological detail

\section{Plain Language Summary}\label{plain-language-summary}

\textbf{In brief}: We compared global ecosystem service models from 1992
and 2020. We identified regions (hotspots) where services changed most
dramatically. We then asked: ``Who lives there?'' (socioeconomic
profiling) and ``Why?'' (land cover drivers). The result is a map of
ecosystem service priorities for conservation and development policy.

\section{Next Steps}\label{next-steps}

Chapter 3 shows what changed globally. Chapters 4-7 drill into WHERE
(geographic patterns), WHO (affected populations), and WHY (land cover
drivers).

\part{Part 2: Global State (WHAT)}

\chapter{Global State (WHAT)}\label{sec-what}

This chapter details the overarching global trajectories of ecosystem
service provision between 1992 and 2020. By aggregating absolute volume
and area changes across major geographic regions (World Bank Regions,
WWF Biomes, and Income Groups), we establish the foundational
understanding of \emph{what} has changed before diving into the
localized extremes (hotspots).

\section{Methodological Note: The ``Aggregate of
Differences''}\label{methodological-note-the-aggregate-of-differences}

To provide the clearest and most accurate narrative of total volume
change, the regional summaries presented in this chapter are derived
from the \textbf{``aggregate of differences''} approach (Path A).

Unlike the 10km grid-based hotspot analysis (Path B), this method
calculates the pixel-by-pixel difference first, and then summarizes
those high-resolution differences directly to the large regional
polygons. This bypasses the Modifiable Areal Unit Problem (MAUP) and
prevents ``Simpson's Paradox'' sign-flips, ensuring that if a region
lost total volume, the chart accurately reflects a negative absolute
change.

\emph{(For a detailed breakdown of the 10km grid-level aggregates and a
discussion of spatial sign-flips, please see Chapter~\ref{sec-annex}).}

\section{Regional Trajectories: Absolute Volume
Change}\label{regional-trajectories-absolute-volume-change}

The following charts display the absolute change in total service
provision across major groupings. For volumetric services (e.g.,
Nitrogen Export, Sediment Export), changes are scaled to the total area
to reflect massive shifts in baseline ecosystems.

\subsection{By World Bank Region}\label{by-world-bank-region}

\begin{figure}

\centering{

\includegraphics[width=1\linewidth,height=\textheight,keepaspectratio]{chapters/../outputs/plots/signed_bars/bars_signed_region_wb_abs.png}

}

\caption{\label{fig-abs-wb}Absolute Change in Ecosystem Services by
World Bank Region}

\end{figure}%

\subsection{By Income Group}\label{by-income-group}

\begin{figure}

\centering{

\includegraphics[width=1\linewidth,height=\textheight,keepaspectratio]{chapters/../outputs/plots/signed_bars/bars_signed_income_grp_abs.png}

}

\caption{\label{fig-abs-income}Absolute Change in Ecosystem Services by
Income Group}

\end{figure}%

\subsection{By WWF Biome}\label{by-wwf-biome}

\begin{figure}

\centering{

\includegraphics[width=1\linewidth,height=\textheight,keepaspectratio]{chapters/../outputs/plots/signed_bars/bars_signed_wwf_biome_abs.png}

}

\caption{\label{fig-abs-biome}Absolute Change in Ecosystem Services by
WWF Biome}

\end{figure}%

\section{Regional Trajectories: Percentage
Change}\label{regional-trajectories-percentage-change}

While absolute change highlights where the largest volumetric losses or
gains occurred (often biased toward massive ecosystems like the Amazon),
percentage change highlights the intensity of the shift relative to the
region's 1992 baseline.

\subsection{By World Bank Region}\label{by-world-bank-region-1}

\begin{figure}

\centering{

\includegraphics[width=1\linewidth,height=\textheight,keepaspectratio]{chapters/../outputs/plots/signed_bars/bars_signed_region_wb_pct.png}

}

\caption{\label{fig-pct-wb}Percentage Change in Ecosystem Services by
World Bank Region}

\end{figure}%

\subsection{By Income Group}\label{by-income-group-1}

\begin{figure}

\centering{

\includegraphics[width=1\linewidth,height=\textheight,keepaspectratio]{chapters/../outputs/plots/signed_bars/bars_signed_income_grp_pct.png}

}

\caption{\label{fig-pct-income}Percentage Change in Ecosystem Services
by Income Group}

\end{figure}%

\subsection{By WWF Biome}\label{by-wwf-biome-1}

\begin{figure}

\centering{

\includegraphics[width=1\linewidth,height=\textheight,keepaspectratio]{chapters/../outputs/plots/signed_bars/bars_signed_wwf_biome_pct.png}

}

\caption{\label{fig-pct-biome}Percentage Change in Ecosystem Services by
WWF Biome}

\end{figure}%

\section{Key Takeaways}\label{key-takeaways-1}

\emph{(This section is reserved for the final interpretation phase. It
will synthesize the core findings from the trajectory charts above.)}

\part{Part 3: Geographic Analysis (WHERE)}

\chapter{Geographic Analysis (WHERE)}\label{sec-where}

Having established the global trajectories in the previous chapter, we
now shift our focus to local extremes. Where is ecosystem service
provision changing the most rapidly relative to local baselines?

We identify \textbf{Hotspots} as the top and bottom 5\% of all 10km grid
cells globally for each service, based on the Symmetric Percentage
Change (SPC) metric. This approach normalizes for the vastly different
baseline sizes of ecosystems across the globe, allowing us to highlight
the most intense local landscape shifts.

\section{The Hotspot Landscape}\label{the-hotspot-landscape}

The following maps provide a ``first look'' at the global distribution
of these extreme changes.

\subsection{First Look: Absolute
Hotspots}\label{first-look-absolute-hotspots}

\emph{(Where the largest total volumes of services were lost or gained)}

\begin{figure}

\centering{

\includegraphics[width=1\linewidth,height=\textheight,keepaspectratio]{chapters/../outputs/plots/maps/first_look_map_abs.png}

}

\caption{\label{fig-map-abs}Global Distribution of Absolute Hotspots}

\end{figure}%

\subsection{First Look: Percentage Hotspots
(Intensity)}\label{first-look-percentage-hotspots-intensity}

\emph{(Where the most intense changes occurred relative to local
baselines)}

\begin{figure}

\centering{

\includegraphics[width=1\linewidth,height=\textheight,keepaspectratio]{chapters/../outputs/plots/maps/first_look_map_pct.png}

}

\caption{\label{fig-map-pct}Global Distribution of Percentage Hotspots}

\end{figure}%

\section{Hotspot Coverage and Relative
Intensity}\label{hotspot-coverage-and-relative-intensity}

How much of each region's land area is experiencing an ecosystem service
hotspot? And is that region disproportionately burdened compared to the
global average?

\subsection{Hotspot Coverage by Income
Group}\label{hotspot-coverage-by-income-group}

\begin{figure}

\centering{

\includegraphics[width=1\linewidth,height=\textheight,keepaspectratio]{chapters/../outputs/plots/intensity/hotspot_coverage_income_grp.png}

}

\caption{\label{fig-cov-income}Area of Hotspots by Income Group}

\end{figure}%

\subsection{Hotspot Relative Intensity by Income
Group}\label{hotspot-relative-intensity-by-income-group}

\emph{(Values \textgreater{} 1 indicate a region has more hotspot area
than expected given its total size)}

\begin{figure}

\centering{

\includegraphics[width=1\linewidth,height=\textheight,keepaspectratio]{chapters/../outputs/plots/intensity/hotspot_relative_intensity_income_grp.png}

}

\caption{\label{fig-rel-income}Relative Intensity of Hotspots by Income
Group}

\end{figure}%

\section{Compound Risk: Multi-Service Hotspots
(``Hotness'')}\label{compound-risk-multi-service-hotspots-hotness}

A single declining service is a concern; multiple co-occurring declines
represent severe compound risk. We refer to the count of overlapping
negative hotspots in a single 10km cell as its ``Hotness''.

\subsection{Compound Risk Map}\label{compound-risk-map}

\begin{figure}

\centering{

\includegraphics[width=1\linewidth,height=\textheight,keepaspectratio]{chapters/../outputs/plots/maps/global_hotspot_count_heatmap_cap4_pct.png}

}

\caption{\label{fig-hotness-map}Global Compound Risk (Hotness) Map}

\end{figure}%

\subsection{Hotness Distribution by World Bank
Region}\label{hotness-distribution-by-world-bank-region}

\begin{figure}

\centering{

\includegraphics[width=1\linewidth,height=\textheight,keepaspectratio]{chapters/../outputs/plots/hotness_dist_region_wb.png}

}

\caption{\label{fig-hotness-wb}Distribution of Overlapping Hotspots by
Region}

\end{figure}%

\subsection{Hotness Distribution by WWF
Biome}\label{hotness-distribution-by-wwf-biome}

\begin{figure}

\centering{

\includegraphics[width=1\linewidth,height=\textheight,keepaspectratio]{chapters/../outputs/plots/hotness_dist_wwf_biome.png}

}

\caption{\label{fig-hotness-biome}Distribution of Overlapping Hotspots
by Biome}

\end{figure}%

\section{Hotspot Distribution Details
(Boxplots)}\label{hotspot-distribution-details-boxplots}

To understand the spread of extreme values within these hotspots, we
look at their distribution across different regions.

\subsection{Volumetric Services by Income
Group}\label{volumetric-services-by-income-group}

\begin{figure}

\centering{

\includegraphics[width=1\linewidth,height=\textheight,keepaspectratio]{chapters/../outputs/plots/boxplots_unified/income_grp/boxplots_volumetric_pct.png}

}

\caption{\label{fig-box-vol}Boxplots of Percentage Hotspots (Volumetric
Services) by Income Group}

\end{figure}%

\subsection{Coastal \& Ratio Services by Income
Group}\label{coastal-ratio-services-by-income-group}

\begin{figure}

\centering{

\includegraphics[width=1\linewidth,height=\textheight,keepaspectratio]{chapters/../outputs/plots/boxplots_unified/income_grp/boxplots_coastal_pct.png}

}

\caption{\label{fig-box-coastal}Boxplots of Percentage Hotspots
(Coastal) by Income Group}

\end{figure}%

\begin{figure}

\centering{

\includegraphics[width=1\linewidth,height=\textheight,keepaspectratio]{chapters/../outputs/plots/boxplots_unified/income_grp/boxplots_ratios_pct.png}

}

\caption{\label{fig-box-ratio}Boxplots of Percentage Hotspots (Ratios)
by Income Group}

\end{figure}%

\section{Summary Statistics}\label{summary-statistics}

The following tables provide the exact numerical breakdown of the
hotspot spatial extraction.

\subsection{Hotspot Area Statistics}\label{hotspot-area-statistics}

{

\begin{longtable}[]{@{}
  >{\raggedright\arraybackslash}p{(\linewidth - 18\tabcolsep) * \real{0.1077}}
  >{\raggedright\arraybackslash}p{(\linewidth - 18\tabcolsep) * \real{0.1769}}
  >{\raggedleft\arraybackslash}p{(\linewidth - 18\tabcolsep) * \real{0.0615}}
  >{\raggedleft\arraybackslash}p{(\linewidth - 18\tabcolsep) * \real{0.0462}}
  >{\raggedleft\arraybackslash}p{(\linewidth - 18\tabcolsep) * \real{0.0692}}
  >{\raggedleft\arraybackslash}p{(\linewidth - 18\tabcolsep) * \real{0.1000}}
  >{\raggedleft\arraybackslash}p{(\linewidth - 18\tabcolsep) * \real{0.0769}}
  >{\raggedright\arraybackslash}p{(\linewidth - 18\tabcolsep) * \real{0.1000}}
  >{\raggedleft\arraybackslash}p{(\linewidth - 18\tabcolsep) * \real{0.1154}}
  >{\raggedleft\arraybackslash}p{(\linewidth - 18\tabcolsep) * \real{0.1462}}@{}}

\caption{\label{tbl-area}Hotspot Area Coverage (Top 10 Rows)}

\tabularnewline

\toprule\noalign{}
\begin{minipage}[b]{\linewidth}\raggedright
service
\end{minipage} & \begin{minipage}[b]{\linewidth}\raggedright
group
\end{minipage} & \begin{minipage}[b]{\linewidth}\raggedleft
n\_total
\end{minipage} & \begin{minipage}[b]{\linewidth}\raggedleft
n\_hot
\end{minipage} & \begin{minipage}[b]{\linewidth}\raggedleft
pct\_area
\end{minipage} & \begin{minipage}[b]{\linewidth}\raggedleft
global\_n\_hot
\end{minipage} & \begin{minipage}[b]{\linewidth}\raggedleft
pct\_share
\end{minipage} & \begin{minipage}[b]{\linewidth}\raggedright
grouping\_var
\end{minipage} & \begin{minipage}[b]{\linewidth}\raggedleft
expected\_share
\end{minipage} & \begin{minipage}[b]{\linewidth}\raggedleft
relative\_intensity
\end{minipage} \\
\midrule\noalign{}
\endhead
\bottomrule\noalign{}
\endlastfoot
Nature\_Access & 4. Lower middle income & 202537 & 16473 & 8.133329 &
67764 & 24.30937 & income\_grp & 15.55465 & 1.562836 \\
Pollination & 4. Lower middle income & 202537 & 15151 & 7.480608 & 67764
& 22.35848 & income\_grp & 15.55465 & 1.437414 \\
N\_export & 4. Lower middle income & 202537 & 13904 & 6.864918 & 67764 &
20.51827 & income\_grp & 15.55465 & 1.319108 \\
N\_Ret\_Ratio & 4. Lower middle income & 202537 & 13606 & 6.717785 &
67764 & 20.07851 & income\_grp & 15.55465 & 1.290836 \\
Sed\_Ret\_Ratio & 4. Lower middle income & 202537 & 13457 & 6.644218 &
67764 & 19.85863 & income\_grp & 15.55465 & 1.276700 \\
Sed\_Ret\_Ratio & 5. Low income & 171032 & 10530 & 6.156743 & 67764 &
15.53922 & income\_grp & 13.13510 & 1.183031 \\
Nature\_Access & 5. Low income & 171032 & 10390 & 6.074887 & 67764 &
15.33263 & income\_grp & 13.13510 & 1.167302 \\
Pollination & 5. Low income & 171032 & 10327 & 6.038051 & 67764 &
15.23966 & income\_grp & 13.13510 & 1.160224 \\
Sed\_export & 4. Lower middle income & 202537 & 11921 & 5.885838 & 67764
& 17.59194 & income\_grp & 15.55465 & 1.130976 \\
Sed\_export & 3. Upper middle income & 603578 & 35028 & 5.803392 & 67764
& 51.69116 & income\_grp & 46.35423 & 1.115134 \\

\end{longtable}

}

\subsection{Hotness Overview}\label{hotness-overview}

{

\begin{longtable}[]{@{}
  >{\raggedright\arraybackslash}p{(\linewidth - 8\tabcolsep) * \real{0.3803}}
  >{\raggedleft\arraybackslash}p{(\linewidth - 8\tabcolsep) * \real{0.1972}}
  >{\raggedleft\arraybackslash}p{(\linewidth - 8\tabcolsep) * \real{0.1408}}
  >{\raggedleft\arraybackslash}p{(\linewidth - 8\tabcolsep) * \real{0.1268}}
  >{\raggedright\arraybackslash}p{(\linewidth - 8\tabcolsep) * \real{0.1549}}@{}}

\caption{\label{tbl-hotness}Compound Risk Summary (Top 10 Rows)}

\tabularnewline

\toprule\noalign{}
\begin{minipage}[b]{\linewidth}\raggedright
group
\end{minipage} & \begin{minipage}[b]{\linewidth}\raggedleft
mean\_hotspots
\end{minipage} & \begin{minipage}[b]{\linewidth}\raggedleft
pct\_multi
\end{minipage} & \begin{minipage}[b]{\linewidth}\raggedleft
n\_pixels
\end{minipage} & \begin{minipage}[b]{\linewidth}\raggedright
grouping
\end{minipage} \\
\midrule\noalign{}
\endhead
\bottomrule\noalign{}
\endlastfoot
4. Lower middle income & 0.4263764 & 9.941393 & 202537 & income\_grp \\
3. Upper middle income & 0.3337647 & 7.616414 & 603578 & income\_grp \\
5. Low income & 0.3284941 & 7.444806 & 171032 & income\_grp \\
1. High income: OECD & 0.2063659 & 4.607142 & 324952 & income\_grp \\
Latin America \& Caribbean & 0.5340240 & 12.791508 & 203768 &
region\_wb \\
East Asia \& Pacific & 0.4522352 & 10.662405 & 245620 & region\_wb \\
Sub-Saharan Africa & 0.3047145 & 6.882040 & 241193 & region\_wb \\
South Asia & 0.2864890 & 5.940008 & 50707 & region\_wb \\
Middle East \& North Africa & 0.2089373 & 5.101566 & 85954 &
region\_wb \\
Europe \& Central Asia & 0.1883946 & 3.915203 & 279679 & region\_wb \\

\end{longtable}

}

\section{Key Takeaways}\label{key-takeaways-2}

\emph{(This section is reserved for the final interpretation phase. It
will synthesize the core findings from the spatial distribution,
intensity, and compound risk charts above.)}

\part{Part 4: Understanding Drivers (WHY)}

\chapter{Understanding Drivers (WHY)}\label{understanding-drivers-why}

\section{The Attribution Gap}\label{the-attribution-gap}

To explain \emph{why} hotspots occur, we integrate Land Cover Change
(LCC) metrics derived from our ESA/C3S data. Instead of simple net
change, we look at the gross loss of natural land to determine how much
of the ecosystem service decline is directly linked to land conversion
versus underlying degradation.

\subsection{Overlap Heatmaps}\label{overlap-heatmaps}

The heatmaps below show what percentage of ecosystem service hotspots
geographically overlap with the top 5\% of land cover conversion (e.g.,
deforestation, urbanization).

\begin{figure}[H]

{\centering \includegraphics[width=1\linewidth,height=\textheight,keepaspectratio]{chapters/../outputs/plots/drivers/heatmap_driver_overlap_pct.png}

}

\caption{Attribution Gap: ES Hotspots vs.~LCC Drivers (Percentage)}

\end{figure}%

\begin{figure}[H]

{\centering \includegraphics[width=1\linewidth,height=\textheight,keepaspectratio]{chapters/../outputs/plots/drivers/heatmap_driver_overlap_abs.png}

}

\caption{Attribution Gap: ES Hotspots vs.~LCC Drivers (Absolute)}

\end{figure}%

\subsection{Correlation Scatterplots}\label{correlation-scatterplots}

The following scatterplots show the relationship between net natural
land change and ecosystem service change.

\begin{figure}[H]

{\centering \includegraphics[width=1\linewidth,height=\textheight,keepaspectratio]{chapters/../outputs/plots/drivers/scatter_Net_Natural_Change_pct.png}

}

\caption{ES Change vs.~Net Natural Change (Percentage)}

\end{figure}%

\begin{figure}[H]

{\centering \includegraphics[width=1\linewidth,height=\textheight,keepaspectratio]{chapters/../outputs/plots/drivers/scatter_Net_Natural_Change_abs.png}

}

\caption{ES Change vs.~Net Natural Change (Absolute)}

\end{figure}%

\part{Part 5: Socioeconomic Analysis (WHO)}

\chapter{Socioeconomic Analysis (WHO)}\label{socioeconomic-analysis-who}

\section{Absolute Population
Exposure}\label{absolute-population-exposure-1}

How many people are directly living within the top 5\% areas of
ecosystem service decline? The charts below break down the exposed
population by World Bank Income Group and the UN Human Development Index
(HDI).

\begin{figure}[H]

{\centering \includegraphics[width=1\linewidth,height=\textheight,keepaspectratio]{chapters/../outputs/plots/pop_exposure_income.png}

}

\caption{Absolute Population Exposure by Income Group}

\end{figure}%

\begin{figure}[H]

{\centering \includegraphics[width=1\linewidth,height=\textheight,keepaspectratio]{chapters/../outputs/plots/pop_exposure_hdi.png}

}

\caption{Absolute Population Exposure by HDI}

\end{figure}%

\section{Relative Socioeconomic
Shift}\label{relative-socioeconomic-shift}

We perform two-sample Kolmogorov-Smirnov (KS) tests to determine if the
socioeconomic profiles of ecosystem service ``hotspots'' are
statistically different from the background landscape.

\begin{figure}[H]

{\centering \includegraphics[width=1\linewidth,height=\textheight,keepaspectratio]{chapters/../outputs/plots/ks/ks_heatmap.png}

}

\caption{KS Statistic Heatmap (Distribution Distance)}

\end{figure}%

\begin{figure}[H]

{\centering \includegraphics[width=1\linewidth,height=\textheight,keepaspectratio]{chapters/../outputs/plots/ks/directionality_cliffs_delta.png}

}

\caption{Directionality of Hotspot Effects (Cliff's Delta)}

\end{figure}%

\chapter{Regional Profiles}\label{regional-profiles}

\emph{(Detailed regional breakdown to be added).}

\part{Part 6: Implications \& Action}

\chapter{Implications \& Action}\label{implications-action}

\emph{(Final conclusions and policy implications to be added).}

\part{Technical Documentation}

\chapter{Methodology: The Two-Path Analysis
Structure}\label{methodology-the-two-path-analysis-structure}

This document clarifies the dual-pathway approach used in the Global NCP
project to analyze changes in ecosystem services between 1992 and 2020.
Understanding this structure is crucial for interpreting the results, as
summary statistics and hotspot analyses are derived from two distinct
computational paths.

\section{The Two Paths}\label{the-two-paths}

The analysis is split into two parallel workflows:

\begin{enumerate}
\def\labelenumi{\arabic{enumi}.}
\tightlist
\item
  \textbf{Path A (Pixel-Level Change):} For raw change summaries.
\item
  \textbf{Path B (Grid-Level Change):} For hotspot identification and
  regional synthesis.
\end{enumerate}

\begin{center}\rule{0.5\linewidth}{0.5pt}\end{center}

\subsection{Path A: Pixel-Level Change
Calculation}\label{path-a-pixel-level-change-calculation}

This path is designed to generate summary statistics (e.g., Mean, Max,
Min) based on the most granular change possible.

\textbf{Workflow:}

\begin{enumerate}
\def\labelenumi{\arabic{enumi}.}
\tightlist
\item
  \textbf{Pixel-wise Difference:} The script
  \texttt{Python\_scripts/batch\_raster\_diff.py} takes the 1992 and
  2020 rasters for a given service and calculates the difference
  (\texttt{T2020\ -\ T1992}) for each corresponding pixel.
\item
  \textbf{Output:} The result is a new set of ``difference rasters''
  that represent the absolute change at the native resolution of the
  data.
\end{enumerate}

\textbf{Use Case:} This path is used when a straightforward,
non-aggregated summary of the overall change across the landscape is
required. \textbf{We use the outputs of this path (via the Python
\texttt{zonal\_stats\_toolkit}) for the main text's regional trajectory
charts (``WHAT'' section) to provide the clearest narrative of total
volume change without spatial aggregation artifacts.}

\subsection{Path B: Grid-Level Change Calculation (The Hotspot
Path)}\label{path-b-grid-level-change-calculation-the-hotspot-path}

This is the primary pathway for the main analysis, including hotspot
identification and regional aggregation. The key distinction is that
\textbf{aggregation precedes differencing}.

\textbf{Workflow:}

\begin{enumerate}
\def\labelenumi{\arabic{enumi}.}
\tightlist
\item
  \textbf{Independent Aggregation:} The script
  \texttt{summary\_pipeline\_landgrid.py} takes the raw 1992 and 2020
  rasters and independently calculates zonal statistics for each of them
  against the canonical 10km grid. The output is a single table
  containing separate columns for the 1992 and 2020 values (e.g.,
  \texttt{service\_1992\_mean}, \texttt{service\_2020\_mean}).
\item
  \textbf{Grid-Level Differencing:} The Quarto notebook
  \texttt{analysis/process\_data.qmd} ingests the aggregated table from
  the previous step. It then calculates the absolute and Symmetric
  Percentage Change (SPC) between the 1992 and 2020 columns for each
  grid cell.
\item
  \textbf{Hotspot Identification:} The notebook
  \texttt{analysis/hotspot\_extraction.qmd} takes the final grid-level
  change data and identifies hotspots.
\end{enumerate}

\textbf{Use Case:} This path underpins all hotspot detection, hotspot
intensity, and socioeconomic exposure analyses (the ``WHERE'', ``WHY'',
and ``WHO''). Its regional summaries of change (the 10km grid averages)
are reserved for the \textbf{Annex} due to the complex aggregation
artifacts described below.

\subsection{Clarification: Rasterizing the Grid vs.~Resampling the
Difference}\label{clarification-rasterizing-the-grid-vs.-resampling-the-difference}

A crucial methodological distinction exists between the output of Path B
and another potential approach.

\begin{itemize}
\item
  \textbf{Rasterizing the Vector Grid (This Project's Method - Path B):}
  We first aggregate the 1992 and 2020 high-resolution data to the 10km
  vector grid, and \emph{then} calculate the difference for each grid
  cell. The final rasters are created by ``burning'' the attribute
  values from this final vector file (\texttt{10k\_change\_calc.gpkg})
  into a GeoTIFF. This is the \textbf{``difference of the aggregates.''}
\item
  \textbf{Resampling the Difference Raster (Path A):} This would involve
  taking the original high-resolution rasters, calculating the
  pixel-by-pixel difference first (creating a ``difference raster''),
  and then aggregating that result to a 10km grid (e.g., by resampling
  or zonal statistics). This would be the \textbf{``aggregate of the
  differences.''}
\end{itemize}

As noted below, these two paths are not always mathematically identical.
This project has officially chosen \textbf{Path B} as the canonical
method for the hotspot analysis. Therefore, the rasterization performed
by \texttt{vector\_to\_raster.py} on the final GeoPackage is the correct
and validated approach for this analysis.

\subsection{\texorpdfstring{Modular Extraction and the \texttt{fid}
Backbone}{Modular Extraction and the fid Backbone}}\label{modular-extraction-and-the-fid-backbone}

A key design feature of Path B is its modularity. The Python extraction
pipeline (\texttt{summary\_pipeline\_landgrid.py}) can be run multiple
times independently (e.g., once for Ecosystem Services, once for
Beneficiaries, once for Coastal Risk). Each run generates a separate
GeoPackage in the workspace.

Because every extraction is performed against the exact same canonical
master grid (\texttt{landgrid\_1\_clean\_enriched\_4326.gpkg}), every
output shares an identical, stable \texttt{fid} (Feature ID) backbone.

To add new variables (e.g., a new year, a new modeled service, or a new
socioeconomic raster) in the future: 1. Run the Python extraction for
the new rasters. 2. The R consolidation script
(\texttt{analysis/process\_data.qmd}) automatically loads the most
recent extraction runs. 3. It performs a rapid, geometry-free tabular
\texttt{left\_join} across all tables using the \texttt{fid}.

This architecture completely prevents the need to re-run heavy,
memory-intensive spatial joins or re-extract existing data when
expanding the analysis.

\begin{center}\rule{0.5\linewidth}{0.5pt}\end{center}

\section{Why Might Results from Path A and Path B
Differ?}\label{why-might-results-from-path-a-and-path-b-differ}

The two paths may produce slightly different summary statistics due to a
fundamental difference in the order of operations:

\begin{itemize}
\tightlist
\item
  \textbf{Path A} calculates the \textbf{aggregate of the differences}.
\item
  \textbf{Path B} calculates the \textbf{difference of the aggregates}.
\end{itemize}

These are not always mathematically identical, especially when dealing
with non-linear metrics or complex spatial distributions. The grid-based
approach (Path B) is considered the canonical method for the hotspot
analysis.

\section{Validation Analysis: Quantifying Methodological Differences
with Path
C}\label{validation-analysis-quantifying-methodological-differences-with-path-c}

To explicitly quantify the difference between the two methodologies and
validate the pipeline, a third analytical path (``Path C'') was
implemented. This path served as a bridge, creating a spatially explicit
grid based on the pixel-level differences.

\textbf{Workflow:}

\begin{enumerate}
\def\labelenumi{\arabic{enumi}.}
\tightlist
\item
  \textbf{Start with Path A Output:} Used the pixel-level ``difference
  rasters'' generated by
  \texttt{Python\_scripts/batch\_raster\_diff.py}.
\item
  \textbf{Aggregate the Difference:} Using a zonal statistics process,
  calculated statistics (e.g., \texttt{mean}, \texttt{sd}) for the pixel
  values from the difference raster within each 10km grid cell.
\end{enumerate}

\textbf{Validation Outcome:}

\begin{itemize}
\tightlist
\item
  \textbf{Aggregate Consistency:} This validation confirmed that the
  results are persistent and mathematically consistent between both
  paths at the aggregate level. The grid of the ``aggregate of the
  differences'' aligns with the Path B grid (``difference of the
  aggregates''), confirming the robustness of the primary grid-level
  hotspot methodology.
\item
  \textbf{Sub-Cell Variability:} By calculating the standard deviation
  (\texttt{sd}) in the aggregation step, we were also able to quantify
  the variability of change \emph{within} each 10km cell. This allows us
  to identify areas of homogenous vs.~heterogeneous change, adding
  valuable richness to the interpretation.
\end{itemize}

This approach provided the spatial representation of the pixel-based
method while remaining computationally manageable, successfully proving
the robustness of the core methodologies before proceeding with the
final hotspot extractions.

\section{Architectural Decision: Spatial Extraction
Strategies}\label{architectural-decision-spatial-extraction-strategies}

Through extensive testing and methodological validation, a distinct
performance and stability divergence was identified between different
zonal statistics engines based on spatial scale and geometry complexity.
The project now officially employs a hybrid extraction strategy:

\begin{enumerate}
\def\labelenumi{\arabic{enumi}.}
\tightlist
\item
  \textbf{For Large Regional Groupings (Complex Multipolygons):
  \texttt{zonal\_stats\_toolkit} (Rasterized)}

  \begin{itemize}
  \tightlist
  \item
    \textbf{Use Case:} Aggregating global data by World Bank Regions,
    Income Groups, or WWF Biomes.
  \item
    \textbf{Rationale:} These groupings form sprawling, highly complex
    multipolygons containing thousands of islands, jagged edges, and
    vertices. Using precise fractional extraction tools
    (\texttt{exactextract} / C++ GEOS) on these shapes causes memory
    leaks and segmentation faults, as the engine attempts to evaluate
    massive geometric intersections. While exploding these polygons into
    tens of thousands of simpler fragments (e.g., 85,000+ pieces)
    prevents crashes, the sheer volume of separate extraction loop
    operations destroys any speed advantage. By rasterizing the polygons
    first (the \texttt{zonal\_stats\_toolkit} approach), we completely
    bypass C++ geometry bottlenecks and reduce the problem to highly
    efficient matrix math.
  \end{itemize}
\item
  \textbf{For High-Resolution Grids (Simple Polygons):
  \texttt{exactextract} (Exact Fractional)}

  \begin{itemize}
  \tightlist
  \item
    \textbf{Use Case:} Aggregating global data into the IUCN 10km
    equal-area grid (millions of cells).
  \item
    \textbf{Rationale:} The 10km grid consists of completely uniform,
    simple square geometries. \texttt{exactextract} is the undisputed
    champion here, chewing through millions of simple shapes almost
    instantly without ever hitting geometry complexity limits.
  \end{itemize}
\end{enumerate}

This ``Best of Both Worlds'' architecture ensures stability and speed
regardless of the spatial target.

\emph{Note: The \texttt{zonal\_stats\_toolkit} is currently managed
alongside the main pipeline in a multi-root VS Code workspace. Future
updates will fully integrate this toolkit's codebase into the
\texttt{global\_NCP} repository to unify version control.}

\section{Spatial Alignment and The Fragment
Bug}\label{spatial-alignment-and-the-fragment-bug}

A major challenge in multi-stage spatial pipelines is maintaining exact
1:1 row integrity between extracted statistics and the canonical master
grid.

\begin{itemize}
\tightlist
\item
  \textbf{The Fragment Bug:} To bypass C++ GEOS bottlenecks during
  Python zonal statistics extraction, complex master grid cells were
  exploded into simpler fragments (e.g., \texttt{gdf.explode()}
  resulting in \textasciitilde1.67M fragments for a 1.5M cell grid). If
  joined directly to downstream datasets, this causes severe data
  duplication (e.g., impossible overlapping hotspot counts per cell) and
  geographic striping.
\item
  \textbf{The Re-aggregation Solution (v1.3.1):} The R pipeline
  (\texttt{process\_data.qmd}) implements a mathematically rigorous
  recovery step. It uses an \texttt{st\_intersects} spatial join
  (mapping fragment centroids to the master grid) to trace every
  fragment back to its pristine 10km parent cell. It then performs a
  re-aggregation (\texttt{group\_by\ \%\textgreater{}\%\ summarise}),
  collapsing the fragments back together. This ensures perfect geometric
  alignment with the canonical grid and strictly bounds the data.
\end{itemize}

\section{Unit Standardization (Per
Hectare)}\label{unit-standardization-per-hectare}

To ensure that comparisons of ecosystem service provision are meaningful
across the globe, all volumetric services (e.g., Nitrogen Export,
Sediment Export) are standardized to a \textbf{per-hectare} basis. This
crucial step corrects for the geometric distortion of raster pixels in
unprojected coordinate systems, where pixel area decreases significantly
with increasing latitude. Without this correction, high-latitude regions
would be disproportionately represented in any analysis based on
per-pixel values.

\begin{itemize}
\tightlist
\item
  \textbf{Volumetric Variables:} Converted to \texttt{Unit\ /\ ha} by
  dividing the raw pixel value by a corresponding pixel area raster
  (\texttt{esa\_pixel\_area\_ha\_...}). This calculation is performed
  \emph{before} the zonal statistics aggregation.
\item
  \textbf{Ratios \& Indices:} Left in their native units (e.g., 0-1
  ratios, unitless indices) as area normalization does not apply.
\end{itemize}

This ensures that the zonal statistics (Mean/Sum) reflect physical
densities comparable across regions.

This per-hectare normalization is a key methodological improvement
(v1.2.1) and is applied to the base year rasters that form the
foundation of \textbf{Path B (Grid-Level Change)}.

\section{Grid Geometry \& Reprojection
Effects}\label{grid-geometry-reprojection-effects}

A common source of confusion in global-scale GIS analysis is the visual
and geometric distortion that occurs during reprojection. This section
clarifies why the project's 10km grid cells may not appear as perfect
squares and confirms the validity of their size.

\begin{itemize}
\tightlist
\item
  \textbf{Geographic vs.~Projected Grids:} The original analysis grid is
  defined in a geographic coordinate system (WGS 84, EPSG:4326), where
  cells are defined by degrees of latitude and longitude. When this grid
  is reprojected into an equal-area system (like Equal Earth, EPSG:8857)
  for analysis, the shapes of the cells are necessarily distorted
  (stretched and tilted) to preserve their area on a flat map.
\item
  \textbf{Bounding Box vs.~True Area:} As a result of this distortion,
  measuring the \texttt{width} and \texttt{height} of a reprojected grid
  cell's bounding box will \textbf{not} yield 10km x 10km. These
  dimensions will vary depending on the cell's latitude.
\item
  \textbf{Area is the Invariant:} The critical property preserved by an
  equal-area projection is the \textbf{area}. A dedicated validation
  script (\texttt{Python\_scripts/verify\_grid\_area.py}) was developed
  to measure the true geometric area of the reprojected grid cells.
\item
  \textbf{Validation Outcome:} This analysis definitively confirmed that
  each grid cell in the reprojected analysis files (e.g.,
  \texttt{...\_epsg8857.gpkg}) has an area of \textbf{100 km²} (10km x
  10km), within a very small tolerance. This validates the integrity of
  the project's spatial foundation and ensures that all area-based
  calculations are correct.
\end{itemize}

\section{Data Preparation: Canonical Master Grid
Generation}\label{data-preparation-canonical-master-grid-generation}

Due to severe performance and geometry-validity bottlenecks encountered
when performing massive spatial joins natively in R (the \texttt{sf}
package), the creation of the canonical 10km master grid was migrated to
a hybrid QGIS-Python workflow.

\begin{enumerate}
\def\labelenumi{\arabic{enumi}.}
\tightlist
\item
  \textbf{Landmass Masking (QGIS):}

  \begin{itemize}
  \tightlist
  \item
    The original global \texttt{AOOGrid} was loaded into QGIS.
  \item
    The WWF Biomes layer was dissolved to create a single global
    landmass polygon.
  \item
    A ``Select by Location'' operation isolated all 10km grid cells
    intersecting the landmass.
  \item
    The selection was exported as \texttt{landgrid\_1.gpkg}.
  \end{itemize}
\item
  \textbf{Geometry Cleaning:}

  \begin{itemize}
  \tightlist
  \item
    To resolve minor topological errors (e.g., self-intersecting rings
    caused by clipping at coastlines and the dateline), the geometries
    were validated and cleaned, resulting in
    \texttt{landgrid\_1\_clean.gpkg}.
  \end{itemize}
\item
  \textbf{Attribute Enrichment (Python):}

  \begin{itemize}
  \tightlist
  \item
    The final stage is handled by
    \texttt{Python\_scripts/enrich\_grid.py}.
  \item
    This script performs a highly optimized, centroid-based spatial join
    to map WWF Biomes, Country, UN/World Bank Regions, and Income Group
    attributes (from the \texttt{ee\_correspondence} dataset) onto the
    grid.
  \item
    It inherently handles geometry validation
    (\texttt{buffer(0).make\_valid()}) and deduplication, outputting the
    final, analysis-ready \texttt{landgrid\_1\_clean\_enriched.gpkg}.
  \end{itemize}
\end{enumerate}

\emph{Note: The legacy R script \texttt{analysis/prepare\_data.qmd} and
standalone cleaning utilities like \texttt{clean\_grid.py} were fully
deprecated in v1.3.4 in favor of this robust Python workflow.}

\section{Data Preparation: Coastal Risk \&
Protection}\label{data-preparation-coastal-risk-protection}

Unlike the other ecosystem services which were natively provided as
continuous global rasters, the Coastal Risk and Protection datasets were
originally provided as high-resolution vector point/line geometries
representing discrete coastal segments.

Attempting vector-on-vector intersections between these intricate
coastal lines and the 1.5-million cell master grid caused severe
performance bottlenecks and \texttt{GEOSException} geometry crashes (due
to the creation of microscopic sliver polygons). To ensure mathematical
robustness and compatibility with the main pipeline, a specialized
pre-processing workflow was utilized:

\begin{enumerate}
\def\labelenumi{\arabic{enumi}.}
\tightlist
\item
  \textbf{Vector Joining \& Ratio Calculation:} The script
  \texttt{Python\_scripts/coastal\_protection\_join.py} merges the 1992
  and 2020 coastal point layers based on their exact spatial locations
  (WKB geometries). It calculates the absolute differences and
  proportional risk reduction ratios (\texttt{Rt\_ratio}) natively in
  the vector domain to prevent floating-point interpolation errors.
\item
  \textbf{Rasterization:} The script
  \texttt{Python\_scripts/rasterize\_coastal.py} safely ``burns'' these
  pre-calculated point values into high-resolution continuous rasters
  (e.g., \texttt{Rt\_1992.tif}, \texttt{Rt\_ratio\_2020.tif}). This uses
  tiled processing to handle the massive global coastline without memory
  exhaustion.
\item
  \textbf{Grid Extraction:} The resulting coastal rasters are seamlessly
  ingested by the standard \texttt{exactextract} pipeline
  (\texttt{summary\_pipeline\_landgrid.py} via
  \texttt{c\_protection\_synth.yaml}). This safely calculates the mean
  coastal risk metrics for each 10km grid cell strictly through
  raster-vector overlap, completely bypassing C-level geometric
  intersection crashes.
\end{enumerate}

By performing the ratio math on the vectors, but the spatial aggregation
on the rasters, we maintain perfect data fidelity without sacrificing
pipeline stability.

\section{Key Analysis Parameters}\label{key-analysis-parameters}

\begin{itemize}
\tightlist
\item
  \textbf{Hotspot Threshold:} The threshold for identifying hotspots is
  defined in \texttt{analysis/hotspot\_extraction.qmd}. It is configured
  in an R object named \texttt{HOTS\_CFG} with the parameter
  \texttt{pct\_cutoff\ =\ 0.05}, representing the top/bottom 5\% of SPC
  values.
\end{itemize}

\textbf{What a ``Hotspot'' Means in This Context:} * \textbf{Relative
Extreme:} We call a cell a ``hotspot'' because its change sits in the
most extreme 5\% \emph{within that specific service's own distribution}.
It is a ranking label, not an absolute physical threshold. A cell can
enter or leave the top 5\% even if its raw change isn't huge in absolute
terms, simply because it is relative to the rest of the globe. *
\textbf{Comparability:} Comparisons are strictly within-service. A top
5\% decline in Service A isn't necessarily comparable in absolute
magnitude to a top 5\% decline in Service B. * \textbf{Not Evidence of
Cause:} Being a hotspot flags that ``this cell's change is unusually
large (relative to peers),'' but it does not inherently prove \emph{why}
the value is extreme. Hotspots are starting points for explanation, not
conclusions. To discuss drivers, we use additional robust analyses (like
LCC attribution and KS socioeconomic profiling). * \textbf{Quick
Analogy:} Top 5\% finishers in two different marathons are both
``elite,'' but their finishing times (absolute performance) and the
reasons they're fast (training, course, weather) can differ. Similarly,
hotspots are elite by rank, not necessarily a proof of shared cause or
identical absolute loss.

\section{Symmetric Percentage Change}\label{symmetric-percentage-change}

To address mathematical artifacts where the sign of percentage change
differs from absolute change (common when baselines are negative or
near-zero), this analysis uses a \textbf{symmetric percentage change}
calculation (\texttt{pct\_mode="symm"}). This ensures that the direction
of the percentage change always aligns with the absolute difference
(\(t_1 - t_0\)).

\textbf{Distribution Limits:} The Symmetric Percentage Change (SPC)
metric is bounded between \textbf{-200\%} (Total Loss) and
\textbf{+200\%} (New Emergence). Consequently, extreme values and
clustering at these boundaries, as well as bi-modal distributions (e.g.,
in Sediment Export), are expected features of the metric rather than
data artifacts.

\subsection{Why SPC over Absolute Change for
Attribution?}\label{why-spc-over-absolute-change-for-attribution}

When assessing the relationship between Land Cover drivers and Ecosystem
Service declines (e.g., scatterplots of Forest Loss vs.~ES Change),
\textbf{Symmetric Percentage Change (SPC)} is strictly preferred over
Absolute Change.

\begin{itemize}
\tightlist
\item
  \textbf{Absolute Change} is inherently biased by the \emph{baseline
  size} of the local ecosystem. A massive, dense forest that loses just
  5\% of its area might show a huge absolute drop in Carbon simply
  because of its initial size. Conversely, a small patch of forest that
  is 100\% destroyed would show a tiny absolute drop. Analyzing Absolute
  Change creates highly skewed, heteroskedastic outputs that largely
  just map ``where the largest baseline ecosystems are.''
\item
  \textbf{Symmetric Percentage Change} normalizes this scale effect. It
  isolates the \emph{intensity of the ecological shock} relative to the
  local baseline. This ensures that a severe multi-service decline in a
  small grid cell is properly recognized as a severe impact, making it
  mathematically appropriate for correlating against land cover
  conversion percentages.
\end{itemize}

\section{Addressing Aggregation Divergence (Simpson's
Paradox)}\label{addressing-aggregation-divergence-simpsons-paradox}

During regional aggregation, it is possible to observe a ``sign
flip''---where regional bar plots display a negative \emph{Absolute
Change} but a positive \emph{Symmetric Percentage Change} (SPC) for the
same service. This is not a calculation error or a bug; it is a
fundamental feature of aggregating spatial data, related to Simpson's
Paradox and the Modifiable Areal Unit Problem (MAUP).

\begin{itemize}
\tightlist
\item
  \textbf{Mean Absolute Change} captures the \textbf{Systemic Shift}.
  Because it averages total physical units, it is heavily weighted by a
  few high-volume grid cells (e.g., massive, dense forests). If those
  highly productive cells lose a large volume of a service, the regional
  absolute average turns negative.
\item
  \textbf{Mean Symmetric Percentage Change} captures the \textbf{Local
  Landscape Shift}. Because percentage change treats every 10km
  community (grid cell) equally regardless of its baseline volume, it
  highlights widespread but low-intensity dynamics.
\end{itemize}

A sign flip reveals a specific geographic narrative: The \emph{total
volume} of the service in the region is decreasing (driven by heavy
localized losses), but the \emph{spatial footprint} of minor expansions
or local service gains is spreading across a large number of
low-baseline cells.

\textbf{Communication Strategy (Main Text vs.~Annex):} Explaining
Simpson's Paradox and spatial aggregation artifacts (MAUP) in the main
text of a paper can severely distract from the core ecological findings.
For this reason: * \textbf{Main Text:} We use \textbf{Path A}
(raster-based regional zonal statistics via the
\texttt{zonal\_stats\_toolkit}) to report the top-level regional
changes. This path calculates total absolute change directly from the
pixels, completely avoiding the MAUP/Simpson's paradox sign-flip issue,
providing a clear and unambiguous ``WHAT'' narrative. * \textbf{Annex:}
The regional trajectory charts derived from \textbf{Path B} (the 10km
grid averages, which exhibit these sign flips) are included in the
Methodological Annex. Path B remains the mathematically robust and
required foundation for detecting local extremes (Hotspots), but its
region-level averages are separated from the main text to preserve
narrative clarity.

\section{Aggregation Logic: Sum
vs.~Mean}\label{aggregation-logic-sum-vs.-mean}

A common question regarding Path B is the comparability of variables
aggregated via \textbf{sum} (extensive variables like Nitrogen Export)
versus those aggregated via \textbf{mean} (intensive variables like Risk
Indices).

\textbf{Why this approach is robust:}

\begin{enumerate}
\def\labelenumi{\arabic{enumi}.}
\tightlist
\item
  \textbf{Physical Correctness:} It is physically correct to sum total
  loads and average representative conditions.
\item
  \textbf{Equal-Area Grid:} The analysis uses the IUCN equal-area grid.
  Since cell area is constant, \(Sum\) and \(Mean\) are perfectly
  proportional (\(Sum = Mean \times Area\)).
\item
  \textbf{Mathematical Identity:} For relative metrics used in this
  analysis (percent change, percentile rankings, KS test statistics),
  the results are identical regardless of whether sum or mean is used.
\item
  \textbf{Comparability:} Comparing relative magnitudes of change (e.g.,
  percentage change) strips away the units, allowing valid comparisons
  between total loads and average indices.
\end{enumerate}

\section{Visualization Semantic Rules
(Maps)}\label{visualization-semantic-rules-maps}

To maintain a consistent narrative across all presentations and figures,
spatial maps of Ecosystem Service change adhere to a strict semantic
color rule:

\begin{itemize}
\tightlist
\item
  \textbf{Universal Diverging Scale:} All maps use a diverging color
  ramp anchored at zero (\texttt{midpoint\ =\ 0}).
\item
  \textbf{Semantic Meaning:} Red always indicates ecological or social
  damage (loss of a good service, or increase in a detrimental risk).
  Green always indicates improvement or healthy service provision.
\item
  \textbf{Sequential vs.~Diverging Data:} Even if a regional dataset
  does not cross zero (e.g., all regions experience a decline), the
  diverging scale is maintained to preserve the semantic meaning of the
  colors. For example, an entirely negative map for a `good' service
  will simply use the Red-to-White half of the scale, avoiding the
  confusion of introducing sequential palettes (like viridis) that break
  the ``Red=Bad'' cognitive mapping.
\end{itemize}

\section{Socioeconomic Profiling (KS
Tests)}\label{socioeconomic-profiling-ks-tests}

To understand the socioeconomic context of ecosystem service hotspots
(e.g., Population, GDP, HDI), we utilize two-sample Kolmogorov-Smirnov
(KS) tests.

\textbf{Balanced Sampling Methodology:} A direct comparison of hotspots
(the top/bottom 5\% of pixels) against the entire non-hotspot background
(the remaining 95\%) suffers from severe sample size imbalance and
includes pixels undergoing extreme changes in the \emph{opposite}
direction. To ensure a fair and stable statistical comparison, the
pipeline implements a ``median background'' sampling strategy: *
Hotspots are compared strictly against the ``business-as-usual'' median
5\% of the landscape (the 47.5th to 52.5th percentiles of change). This
isolates the specific socioeconomic profile of extreme decline against
typical, stable baseline conditions.

\section{Land Cover Change
Attribution}\label{land-cover-change-attribution}

To explain \emph{why} hotspots occur, we integrate Land Cover Change
(LCC) metrics derived from ESA CCI (1992) and C3S (2020) maps.

\textbf{Methodology:} Instead of simple ``Net Change'' (which masks
simultaneous loss and gain), we use the \textbf{\texttt{diffeR}
methodology} (Pontius et al.) to calculate: * \textbf{Gross Loss:} The
specific area of natural land lost to transformation. * \textbf{Gross
Gain:} The area of natural land recovered. * \textbf{Exchange:} Shifts
that don't affect the net total but represent dynamic turnover.

These metrics are aggregated to the 10km grid and overlaid with ES
hotspots to quantify the \textbf{``Attribution Gap''} (i.e., how much ES
decline is directly linked to land conversion vs.~degradation). The
analysis produces both a single, global attribution map showing the
overall footprint of degradation, as well as a series of detailed maps
breaking down the specific drivers for each of the 8 individual
ecosystem service hotspots.

\textbf{Granular Models:} To move beyond binary ``Natural
vs.~Transformed'' analysis, we implement two specific driver models:

\begin{enumerate}
\def\labelenumi{\arabic{enumi}.}
\tightlist
\item
  \textbf{Forest Loss Model:}

  \begin{itemize}
  \tightlist
  \item
    \textbf{Reclassification:} Maps ESA classes to \textbf{Forest}
    vs.~\textbf{Non-Forest}. Crucially, Flooded Trees (classes 160, 170)
    are mapped to Forest to capture mangrove/swamp forest dynamics.
  \item
    \textbf{Metric:} Tracks Gross Loss of Forest cover.
  \end{itemize}
\item
  \textbf{Expansion Model:}

  \begin{itemize}
  \tightlist
  \item
    \textbf{Reclassification:} Maps ESA classes to \textbf{Urban},
    \textbf{Cropland}, and \textbf{Other}.
  \item
    \textbf{Metric:} Tracks the specific expansion of Urban and Cropland
    areas into other land cover types.
  \end{itemize}
\end{enumerate}

\textbf{Note on Rangelands:} Logic is updated to explicitly track
Forest-to-Grassland transitions as a loss of primary natural cover.
Categorizing Grasslands/Shrublands as
\texttt{Transformed\ (Rangeland/Pasture)} prevents these critical
conversions from being masked as `Natural-to-Natural' exchange.

\chapter{Methodological Annex: 10km Grid Aggregation}\label{sec-annex}

This annex provides supplemental detail on the primary spatial framework
underpinning the hotspot analysis in this book: \textbf{The 10km
Equal-Area Grid (Path B)}.

\section{The Two-Path Analysis
Structure}\label{the-two-path-analysis-structure}

As detailed in the full technical methodology
(\texttt{docs/methodology.md}), this project utilizes a dual-pathway
approach to handle the complexities of global spatial data.

\begin{itemize}
\tightlist
\item
  \textbf{Path A (Aggregate of Differences):} This path, presented in
  Chapter~\ref{sec-what}, calculates pixel-by-pixel change first, and
  then sums those changes to large regional polygons. It is used to
  present the clearest picture of total volumetric shifts without
  spatial aggregation artifacts.
\item
  \textbf{Path B (Difference of Aggregates):} This is the canonical path
  for all hotspot detection, intensity mapping, attribution, and
  socioeconomic profiling (Chapter~\ref{sec-where}, \textbf{?@sec-why},
  \textbf{?@sec-who}). It aggregates the raw 1992 and 2020 data to the
  10km grid first, and then calculates the difference for each grid
  cell.
\end{itemize}

\section{The Modifiable Areal Unit Problem (MAUP) and
``Sign-Flips''}\label{the-modifiable-areal-unit-problem-maup-and-sign-flips}

When we average the cell-level changes from Path B up to regional
scales, we sometimes observe a phenomenon related to Simpson's Paradox
and the Modifiable Areal Unit Problem (MAUP): \textbf{Sign-Flips}.

A ``sign-flip'' occurs when a region shows a negative \emph{Absolute
Change} but a positive \emph{Symmetric Percentage Change} (SPC) for the
same service.

This is not an error; it reveals a specific geographic narrative: *
\textbf{Mean Absolute Change} is heavily weighted by a few high-volume
grid cells (e.g., massive, dense forests). If those highly productive
cells lose a large volume of a service, the regional absolute average
turns negative. * \textbf{Mean Symmetric Percentage Change} treats every
10km community equally, regardless of its baseline volume. It highlights
widespread, low-intensity dynamics.

A sign-flip indicates that the \emph{total volume} of the service in the
region is decreasing (driven by heavy localized losses), but the
\emph{spatial footprint} of minor expansions or local service gains is
spreading across a large number of low-baseline cells.

\section{10km Grid Regional
Aggregations}\label{km-grid-regional-aggregations}

The following charts display the regional summaries calculated via Path
B. Compare these with the Path A charts in Chapter~\ref{sec-what} to
observe the subtle differences introduced by the 10km spatial
aggregation framework.

\subsection{By World Bank Region (10km
Aggregation)}\label{by-world-bank-region-10km-aggregation}

\begin{figure}

\centering{

\includegraphics[width=1\linewidth,height=\textheight,keepaspectratio]{chapters/../outputs/plots/signed_bars/bars_signed_region_wb_abs.png}

}

\caption{\label{fig-annex-abs-wb}10km Grid Aggregation: Absolute Change
by Region}

\end{figure}%

\begin{figure}

\centering{

\includegraphics[width=1\linewidth,height=\textheight,keepaspectratio]{chapters/../outputs/plots/signed_bars/bars_signed_region_wb_pct.png}

}

\caption{\label{fig-annex-pct-wb}10km Grid Aggregation: Percentage
Change by Region}

\end{figure}%

\subsection{By Income Group (10km
Aggregation)}\label{by-income-group-10km-aggregation}

\begin{figure}

\centering{

\includegraphics[width=1\linewidth,height=\textheight,keepaspectratio]{chapters/../outputs/plots/signed_bars/bars_signed_income_grp_abs.png}

}

\caption{\label{fig-annex-abs-inc}10km Grid Aggregation: Absolute Change
by Income Group}

\end{figure}%

\begin{figure}

\centering{

\includegraphics[width=1\linewidth,height=\textheight,keepaspectratio]{chapters/../outputs/plots/signed_bars/bars_signed_income_grp_pct.png}

}

\caption{\label{fig-annex-pct-inc}10km Grid Aggregation: Percentage
Change by Income Group}

\end{figure}%

\subsection{By WWF Biome (10km
Aggregation)}\label{by-wwf-biome-10km-aggregation}

\begin{figure}

\centering{

\includegraphics[width=1\linewidth,height=\textheight,keepaspectratio]{chapters/../outputs/plots/signed_bars/bars_signed_wwf_biome_abs.png}

}

\caption{\label{fig-annex-abs-biome}10km Grid Aggregation: Absolute
Change by Biome}

\end{figure}%

\begin{figure}

\centering{

\includegraphics[width=1\linewidth,height=\textheight,keepaspectratio]{chapters/../outputs/plots/signed_bars/bars_signed_wwf_biome_pct.png}

}

\caption{\label{fig-annex-pct-biome}10km Grid Aggregation: Percentage
Change by Biome}

\end{figure}%




\end{document}
